\section*{Original main statement (journal p.7)}
For $n,d\in\mathbb N^+$ and $k\in\{0,1,\ldots,n-1\}$, define $N(k,d,n)$ by
\[N(k,d,n):=\min\left\{\binom{(k'+1)(n'+1+d)}{(k'+1)d},\binom{(n-k)(n'+1+d)}{(n-k)d}\right\},\]
where $k':=\min\{k,d-2\}\le k$ and $n':=k'+n-k\le n$.
\begin{result}{Theorem 1.7 (Main Theorem)}
For $n,d\in\mathbb N^+$, $k\in\{0,1,\ldots,n-1\}$, and $\epsilon\in(0,1)$, there exists an $(n,d,\epsilon)$-evasive $k$-subspace family (resp. affine $k$-subspace family) $\mathcal H$ on $\PP^n$ (resp. $\Aspace^n$) of size
\[\operatorname{poly}(N(k,d,n),n,1/\epsilon)\le\operatorname{poly}(n^{\min\{k+1,n-k,d\}d},d,1/\epsilon).\]
Moreover:
\begin{itemize}
\item If $\operatorname{char}(\FF)=0$, where $\FF$ denotes the base field, then the linear equations defining the projective or affine subspaces in $\mathcal H$ are defined over $\mathbb Q$. Moreover, the bit-lengths of the numerators and denominators of the coefficients of these linear equations are polynomial in $|\mathcal H|$, and the total bit complexity of computing these linear equations is polynomial in $|\mathcal H|$.
\item If $\operatorname{char}(\FF)=p>0$, then the linear equations defining the projective or affine subspaces in $\mathcal H$ are defined over an extension field $\FF_q$ of $\FF_p$, where $q\le\operatorname{poly}(|\mathcal H|,p)$. The total bit complexity of computing these linear equations and constructing the field $\FF_q$ is polynomial in $|\mathcal H|$ and $\log p$.
\end{itemize}
In particular, when $d$ is bounded, the bit complexity of constructing $\mathcal H$ is polynomial in $n/\epsilon$ (and $\log p$ if $\operatorname{char}(\FF)=p>0$).
\end{result}
